\chapter{Anatomie d'un prédicteur}\label{ch:rs:anatomy-of-a-predictor}

Deux chercheurs rapportent le \omterm{def:rs:anatomy-of-a-predictor:ic}{coefficient d'information} du même signal de retournement sur les mêmes actions~: 0.038 et 0.007. Les deux
ont calculé une corrélation de rang entre l'opposé d'un rendement passé et un rendement futur, et tous deux l'ont fait correctement. Le premier
a corrélé le rendement d'hier avec celui de demain~; le second, le rendement des cinq derniers jours avec celui des cinq jours suivants. Un \omterm{def:rs:anatomy-of-a-predictor:predictor}{prédicteur}
n'est pas une formule~: c'est une formule, une cible, un horizon, une normalisation et une \omterm{def:rs:anatomy-of-a-predictor:neutralise}{neutralisation}, et un nombre qui le décrit ne signifie
rien sans ces cinq éléments. Ce chapitre décompose un \omterm{def:rs:anatomy-of-a-predictor:predictor}{prédicteur} en ces pièces et le mesure avec le \omterm{def:rs:anatomy-of-a-predictor:ic}{coefficient d'information}
et les erreurs types que la mesure exige, et se termine par la \omterm{def:rs:anatomy-of-a-predictor:card}{fiche de prédicteur}, la fiche d'une page que
le reste de la partie~II remplit pour chaque \omterm{def:rs:anatomy-of-a-predictor:predictor}{prédicteur} étudié. Ses données sont le marché synthétique du livre (chapitre~5), dont
le retournement d'un jour est connu pour être réel.

\section{La cible}

\begin{definition}[Prédicteur, cible de prévision, horizon de prévision]\label{def:rs:anatomy-of-a-predictor:predictor}
Un \emph{prédicteur}\index{prédicteur} est un nombre calculé pour chaque titre à un instant de décision à partir de données connues à cet instant,
destiné à être corrélé, d'un titre à l'autre, à une quantité ultérieure, sa \emph{cible de prévision}\index{cible de prévision}
(en général un rendement). L'\emph{horizon de prévision}\index{horizon de prévision} est la durée de la cible~: le rendement sur le
jour suivant, sur les cinq jours suivants, sur le mois suivant.
\end{definition}

\begin{omfigure}
\begin{tikzpicture}[font=\footnotesize,x=1.1cm,y=1cm]
\draw[->,>=Stealth] (0,0) -- (11.2,0) node[right] {jours de négociation};
\foreach \x in {0,1,...,10} \draw (\x,0.08) -- (\x,-0.08);
\node[below right] at (5.05,-0.05) {$t$};
\fill[omDef!20] (0,0.25) rectangle (4.9,0.75);
\node[omDef] at (2.45,0.5) {fenêtre du prédicteur~: données jusqu'à $t$};
\fill[omThm!20] (5,-0.75) rectangle (10,-1.25);
\node[omThm] at (7.5,-1.0) {cible~: rendement sur $t+1, \dots, t+h$};
\draw[very thick,black] (5,1.0) -- (5,-1.5);
\node[above] at (5,1.0) {décision à la clôture de $t$};
\end{tikzpicture}

\omcaption{Le calendrier d'un \omterm{def:rs:anatomy-of-a-predictor:predictor}{prédicteur}. Tout ce qu'il utilise est connu à la clôture du jour de décision $t$~; sa cible commence
le jour suivant. Un \omterm{def:rs:anatomy-of-a-predictor:predictor}{prédicteur} dont la fenêtre dépasse $t$, ou une cible qui commence en $t$, comporte une anticipation (chapitre~3).}\label{fig:rs:anatomy-of-a-predictor:timing}
\end{omfigure}

L'horizon fait partie du \omterm{def:rs:anatomy-of-a-predictor:predictor}{prédicteur}. Sur le marché synthétique, le retournement mesuré de quatre façons donne quatre nombres différents~:
l'opposé du rendement du dernier jour a un \omterm{def:rs:anatomy-of-a-predictor:ic}{coefficient d'information de rang} moyen de 0.0385 avec le rendement du jour suivant ($t = 17.8$), et de 0.0162 avec le
rendement des cinq jours suivants~; l'opposé du rendement des cinq derniers jours en a 0.0187 avec le jour suivant et 0.0074 avec les cinq jours suivants ($t = 2.4$).
Le simulateur introduit un retournement d'un jour, et toute mesure qui dilue ce jour, de part et d'autre, en voit moins. Pour un
\omterm{def:rs:anatomy-of-a-predictor:predictor}{prédicteur} réel, l'horizon est inconnu et doit être trouvé (chapitre~13), mais la leçon demeure~: le couple (fenêtre du \omterm{def:rs:anatomy-of-a-predictor:predictor}{prédicteur},
horizon de la cible) qui est mesuré.

\begin{definition}[Rendement résiduel]\label{def:rs:anatomy-of-a-predictor:residual}
Le \emph{rendement résiduel}\index{rendement résiduel} d'un titre sur une période est son rendement diminué de la part expliquée par
des expositions aux facteurs choisies, estimée à chaque date par une régression transversale des rendements sur ces expositions~: ce qui
reste après le marché, les secteurs ou les styles.
\end{definition}

Trois cibles sont courantes~: le rendement brut, le rendement en excès d'un marché (ou d'un indice de référence), et un \omterm{def:rs:anatomy-of-a-predictor:residual}{rendement résiduel}. Pour un \omterm{def:rs:anatomy-of-a-predictor:ic}{coefficient d'information de rang}
calculé entre titres à une date donnée, les deux premières sont identiques, puisque retrancher un même nombre à chaque rendement laisse
leurs rangs inchangés~: sur le marché synthétique, elles donnent toutes deux 0.0074 pour le retournement sur cinq jours. Le résiduel est différent, et la
section suivante précise quand il compte.

\section{Normalisation}

\begin{definition}[Z-score transversal, transformation en rangs]\label{def:rs:anatomy-of-a-predictor:normalise}
Le \emph{z-score transversal}\index{z-score transversal} d'un \omterm{def:rs:anatomy-of-a-predictor:predictor}{prédicteur} à une date est sa valeur diminuée de la moyenne de cette date
entre titres, divisée par l'écart-type de cette date, généralement tronqué à quelques unités. Sa \emph{transformation en
rangs}\index{transformation en rangs} remplace chaque valeur par son rang à cette date, ramené dans l'intervalle $(-\tfrac12, \tfrac12)$.
\end{definition}

Les \omterm{def:rs:anatomy-of-a-predictor:predictor}{prédicteurs} bruts ne sont comparables ni d'une date à l'autre (un rendement de 3~\% sur cinq jours est important lors d'une semaine calme et faible lors d'un krach)
ni d'un \omterm{def:rs:anatomy-of-a-predictor:predictor}{prédicteur} à l'autre (un rendement et un ratio valeur comptable sur cours ont des unités différentes). Normaliser la coupe transversale de chaque date
supprime ces deux difficultés et transforme un \omterm{def:rs:anatomy-of-a-predictor:predictor}{prédicteur} en un ensemble de paris relatifs, ce que détient un livre market-neutral. Le
z-score conserve la forme de la distribution et doit être tronqué (ou le \omterm{def:rs:anatomy-of-a-predictor:predictor}{prédicteur} winsorisé, livre~4, chapitre~15), car un
seul titre extrême dominerait sinon la date. La \omterm{def:rs:anatomy-of-a-predictor:normalise}{transformation en rangs} élimine entièrement la forme~: elle est robuste, et le \omterm{def:rs:anatomy-of-a-predictor:ic}{coefficient d'information de rang},
la corrélation de Spearman du livre~4, chapitre~15, est la corrélation de Pearson de deux variables transformées en rangs. Rien n'est gratuit~:
la \omterm{def:rs:anatomy-of-a-predictor:normalise}{transformation en rangs} traite l'écart entre le premier et le deuxième titre comme celui entre le 500e et le 501e.

\section{Neutralisation}

\begin{definition}[Neutralisation]\label{def:rs:anatomy-of-a-predictor:neutralise}
La \emph{neutralisation}\index{neutralisation} consiste à retirer d'un \omterm{def:rs:anatomy-of-a-predictor:predictor}{prédicteur}, ou d'une cible, sa composante le long d'expositions choisies
en prenant à chaque date le résidu d'une régression transversale sur ces expositions~: la neutralisation sectorielle retire
les moyennes sectorielles, la neutralisation au bêta retire la composante le long du bêta de marché.
\end{definition}

\begin{proposition}[Ce que la neutralisation de la cible fait au coefficient d'information]\label{prop:rs:anatomy-of-a-predictor:neutral}
À une date donnée, écrivons la cible sous la forme $y = f + u$, où $f$ est sa projection sur les expositions et $u$ le résidu, non corrélé
à $f$. Si le \omterm{def:rs:anatomy-of-a-predictor:predictor}{prédicteur} $s$ est non corrélé à $f$, alors
\[
\operatorname{corr}(s, u) = \operatorname{corr}(s, y)\,\sqrt{\frac{\Var(y)}{\Var(u)}}.
\]
\end{proposition}

\begin{proof}
$\Cov(s, y) = \Cov(s, u)$, puisque $s$ est non corrélé à $f$. Diviser par $\sqrt{\Var(s)\Var(y)}$ et par $\sqrt{\Var(s)\Var(u)}$
tour à tour.
\end{proof}

Si les expositions expliquent une part $1 - \rho$ de la variance transversale de la cible, mesurer par rapport au résiduel élève le
\omterm{def:rs:anatomy-of-a-predictor:ic}{coefficient d'information} de $1/\sqrt\rho$~: le \omterm{def:rs:anatomy-of-a-predictor:predictor}{prédicteur} avait raison sur la partie qu'il voyait, et la composante factorielle n'était pour lui que du bruit. Dans le
marché synthétique par défaut, le bêta de marché et dix secteurs n'expliquent que 10~\% de la variance des rendements sur cinq jours entre actions, et
la \omterm{def:rs:anatomy-of-a-predictor:neutralise}{neutralisation} fait passer le \omterm{def:rs:anatomy-of-a-predictor:ic}{coefficient d'information} du retournement sur cinq jours de 0.0074 à 0.0064, dans le bruit. Dans une variante dont les secteurs fluctuent davantage
(30~\% de volatilité annuelle au lieu de 8~\%), ils en expliquent 52~\%, et la proposition prédit un rapport de $1/\sqrt{0.476} = 1.45$~: le \omterm{def:rs:anatomy-of-a-predictor:ic}{coefficient d'information}
passe de 0.0043 à 0.0067 par rapport à la cible résiduelle, et à 0.0086 lorsque le \omterm{def:rs:anatomy-of-a-predictor:predictor}{prédicteur} est lui aussi neutralisé par secteur, la
statistique $t$ passant de 0.8 à 3.5 (\cref{fig:rs:anatomy-of-a-predictor:neutral}). Neutraliser le \omterm{def:rs:anatomy-of-a-predictor:predictor}{prédicteur} retire sa propre
composante sectorielle, laquelle ne porte dans ce marché aucune information sur le futur.

\begin{omfigure}
\begin{tikzpicture}
\begin{axis}[width=9cm,height=5cm,ybar,bar width=26pt,ylabel={coefficient d'information de rang moyen},tick label style={font=\small},
  label style={font=\small},xmin=-0.6,xmax=2.6,ymin=0,ymax=0.011,xtick={0,1,2},
  xticklabels={{brut},{cible\\résiduelle},{les deux\\neutralisés}},xticklabel style={align=center},scaled y ticks=false,
  yticklabel style={/pgf/number format/fixed,/pgf/number format/precision=3},grid=major,grid style={gray!25},
  nodes near coords,point meta=explicit symbolic,every node near coord/.append style={font=\footnotesize},
  table/col sep=comma]
\addplot[fill=omDef!55,draw=omDef] table[x=k,y=ic,meta=t]{figdata/research/06-anatomy-of-a-predictor/neutral.csv};
\end{axis}
\end{tikzpicture}

\omcaption{Retournement à cinq jours contre rendements à cinq jours dans un marché synthétique dont les secteurs bougent fortement (30~\% par an)~: \omterm{def:rs:anatomy-of-a-predictor:predictor}{prédicteur}
brut et cible brute, \omterm{def:rs:anatomy-of-a-predictor:predictor}{prédicteur} brut et cible résiduelle, \omterm{def:rs:anatomy-of-a-predictor:predictor}{prédicteur} neutralisé par secteur et cible résiduelle. Les étiquettes sont
les statistiques $t$ HAC. Données~: \texttt{firm.synthmkt} avec \texttt{ind\_vol=0.30}, graine 1.}\label{fig:rs:anatomy-of-a-predictor:neutral}
\end{omfigure}

\section{Le coefficient d'information et son erreur type}

\begin{definition}[Coefficient d'information, coefficient d'information de rang]\label{def:rs:anatomy-of-a-predictor:ic}
Le \emph{coefficient d'information}\index{coefficient d'information} (IC) d'un \omterm{def:rs:anatomy-of-a-predictor:predictor}{prédicteur} à une date donnée est la corrélation, sur
les titres de cette date, entre le \omterm{def:rs:anatomy-of-a-predictor:predictor}{prédicteur} et sa cible. Le \emph{coefficient d'information de rang}%
\index{coefficient d'information de rang} est le même calcul effectué sur les rangs (corrélation de Spearman). Un \omterm{def:rs:anatomy-of-a-predictor:predictor}{prédicteur} se résume par la série temporelle de ses
IC~: leur moyenne, leur écart type et le rapport des deux.
\end{definition}

\begin{definition}[Ratio d'information de l'IC, écart entre quantiles]\label{def:rs:anatomy-of-a-predictor:icir}
Le \emph{ratio d'information de l'IC}\index{ratio d'information de l'IC} (ICIR) est l'IC moyen divisé par son écart type sur les
dates, annualisé par $\sqrt{252/h}$ pour un horizon de $h$ jours échantillonné tous les $h$ jours. L'\emph{écart entre quantiles}\index{ecart entre quantiles@écart entre quantiles} est
la cible moyenne des titres du quantile supérieur du \omterm{def:rs:anatomy-of-a-predictor:predictor}{prédicteur} moins celle du quantile inférieur, date par date.
\end{definition}

Deux faits sur le bruit de l'IC déterminent la durée pendant laquelle un \omterm{def:rs:anatomy-of-a-predictor:predictor}{prédicteur} doit être observé. À une date donnée, avec $N$ titres
et aucune prévisibilité, l'IC a un écart type d'environ $1/\sqrt{N - 1}$ (0.032 pour 1\,000 titres)~; une moyenne d'IC de 0.01 est donc
invisible à une date isolée et doit être accumulée sur des centaines de dates. Et lorsque la cible couvre $h$ jours mais que l'IC est calculé chaque jour, les
cibles de dates consécutives partagent $h - 1$ jours, les IC sont fortement autocorrélés, et
l'erreur type naïve de leur moyenne est trop petite.

\begin{proposition}[Cibles qui se chevauchent]\label{prop:rs:anatomy-of-a-predictor:overlap}
Si la série quotidienne des IC d'une cible à $h$ jours a des autocorrélations $\rho_k$ qui décroissent linéairement jusqu'à zéro au retard $h$ (comme pour un
\omterm{def:rs:anatomy-of-a-predictor:predictor}{prédicteur} persistant et des rendements quotidiens non chevauchants), la variance de sa moyenne sur $n$ dates vaut approximativement $h$ fois la
variance naïve $\Var/n$~: la statistique $t$ naïve surestime jusqu'à $\sqrt h$. L'estimateur de Newey--West à $h - 1$ retards (livre~4,
chapitre~11) la corrige, de même que le fait de ne retenir qu'une date sur $h$.
\end{proposition}

\begin{proof}
$\Var(\bar x) \approx \frac{\sigma^2}{n}\bigl(1 + 2\sum_{k=1}^{h-1}\rho_k\bigr)$ avec $\rho_k = 1 - k/h$~; la somme vaut $(h - 1)/2$, donc le
crochet vaut $h$.
\end{proof}

Pour le retournement à cinq jours contre des cibles à cinq jours, la statistique $t$ naïve vaut 3.85 et la statistique corrigée 2.39~: le rapport 1.61 est
inférieur à $\sqrt5 = 2.24$ parce que le retournement n'est pas persistant. Ne retenir qu'une date sur cinq donne 1.31, une estimation valide mais moins
efficace, fondée sur un cinquième des données. Pour le retournement à un jour contre des rendements à un jour, il n'y a pas de chevauchement~: 0.0385, $t = 17.8$, un
ICIR annualisé de 6.4 une fois le \omterm{def:rs:anatomy-of-a-predictor:predictor}{prédicteur} neutralisé, un chiffre qu'aucun \omterm{def:rs:anatomy-of-a-predictor:predictor}{prédicteur} quotidien réel n'atteint après coûts~; le simulateur
l'inscrit fortement pour que les outils des chapitres suivants aient quelque chose à trouver.

\begin{omfigure}
\begin{tikzpicture}
\begin{axis}[width=5.8cm,height=4.8cm,ybar,bar width=7pt,xlabel={jour $k$ après la décision},ylabel={IC de rang avec le rendement du jour $k$},
  tick label style={font=\small},label style={font=\small},xmin=0.3,xmax=10.7,ymin=-0.01,ymax=0.045,xtick={1,2,3,4,5,6,7,8,9,10},
  scaled y ticks=false,yticklabel style={/pgf/number format/fixed,/pgf/number format/precision=2},grid=major,grid style={gray!25},
  table/col sep=comma]
\addplot[fill=omDef!55,draw=omDef] table[x=k,y=ic]{figdata/research/06-anatomy-of-a-predictor/lags.csv};
\end{axis}
\end{tikzpicture}\hspace{3mm}%
\begin{tikzpicture}
\begin{axis}[width=5.8cm,height=4.8cm,xlabel={années},ylabel={IC cumulé},tick label style={font=\small},
  label style={font=\small},xmin=0,xmax=10,ymin=0,grid=major,grid style={gray!25},table/col sep=comma]
\addplot[omDef,very thick] table[x=year,y=cum]{figdata/research/06-anatomy-of-a-predictor/cumic.csv};
\end{axis}
\end{tikzpicture}

\omcaption{Le retournement à un jour sur le marché synthétique. À gauche~: son IC de rang avec le rendement de chacun des dix jours suivants~; la
information est épuisée dès le premier jour (0.0385), le second n'en porte que 0.0008. À droite~: la somme cumulée de son IC quotidien sur dix ans,
neutralisé au bêta et aux secteurs~; une droite est le signe d'un \omterm{def:rs:anatomy-of-a-predictor:predictor}{prédicteur} stable. Données~: \texttt{firm.synthmkt}, graine 1.}\label{fig:rs:anatomy-of-a-predictor:lags}
\end{omfigure}

\section{La fiche de prédicteur}

\begin{definition}[Fiche de prédicteur]\label{def:rs:anatomy-of-a-predictor:card}
Une \emph{fiche de prédicteur}\index{fiche de prédicteur} est le relevé d'une page qui décrit un \omterm{def:rs:anatomy-of-a-predictor:predictor}{prédicteur}~: sa définition sous forme de formule~; ses entrées et
leur \omterm{def:rs:point-in-time-data-and-the-biases:bitemporal}{temps de connaissance}~; sa justification~; son horizon et sa demi-vie mesurée~; sa normalisation et sa \omterm{def:rs:anatomy-of-a-predictor:neutralise}{neutralisation}~; ses
modes de défaillance connus~; ses sources. Les statistiques mesurées proviennent d'un script testé, jamais de la mémoire.
\end{definition}

La fiche est l'unité dans laquelle la partie~II de ce livre, et un groupe de recherche, accumulent le savoir~: un \omterm{def:rs:anatomy-of-a-predictor:predictor}{prédicteur} sans fiche
n'a pas été étudié, seulement essayé. Ses champs obligent à répondre aux questions de ce chapitre avant que le \omterm{def:rs:anatomy-of-a-predictor:predictor}{prédicteur} ne soit utilisé. Voici
la première.

\begin{predictorcard}{Retournement à un jour}\label{pred:rs:anatomy-of-a-predictor:reversal}
\sfield{Définition} $s_{i,t} = -r_{i,t}$, le rendement du jour au signe changé, classé à travers l'univers et neutralisé au bêta de
marché et aux secteurs.
\sfield{Entrées et horodatages} Les clôtures des jours $t-1$ et $t$, connues à la clôture de $t$~; bêtas et secteurs à la date $t$.
\sfield{Justification} Les fournisseurs de liquidité sont rémunérés, par un retournement, pour avoir absorbé la demande d'immédiateté (chapitre~1).
\sfield{Horizon et demi-vie} Un jour. Sur le marché synthétique, l'IC de rang vaut 0.0385 avec le jour $t+1$ et 0.0008 avec le jour
$t+2$~: la demi-vie est inférieure à un jour.
\sfield{Normalisation} Rang transversal~; résidu sur le bêta et les indicatrices sectorielles. Neutralisé, son IC de rang moyen vaut
0.040 avec un ICIR annualisé de 6.4 ($t = 20.1$) sur le marché synthétique, qui l'inscrit.
\sfield{Modes de défaillance} Les mouvements liés aux annonces (jours de résultats) ne reviennent pas~; la rotation quotidienne coûte
plus que ce que rapporte l'IC dans les grandes valeurs liquides~; l'encombrement (crowding) dû à de nombreux fournisseurs de liquidité.
\sfield{Sources} N.~Jegadeesh (1990)~; B.~N.~Lehmann (1990)~; \texttt{rs\_anatomy.card()} pour les chiffres.
\end{predictorcard}

\section{Tutoriel~: mesurer un prédicteur}\label{tut:rs:anatomy-of-a-predictor:measure}

\begin{tutorial}
\textbf{Objectif.} Mesurer le retournement sur le marché synthétique avec chacun des choix de ce chapitre rendus explicites, et remplir sa fiche.
\textbf{État final~:} les \cref{fig:rs:anatomy-of-a-predictor:neutral,fig:rs:anatomy-of-a-predictor:lags}~; le tableau des horizons~; les
statistiques $t$ 3.85 (naïve) et 2.39 (corrigée)~; la \cref{pred:rs:anatomy-of-a-predictor:reversal}.
\begin{tutsteps}
\item \textbf{Neutraliser.} Une régression transversale par date sur une constante et les expositions~; les titres dont une valeur manque
sont écartés, non complétés.
\omcode{code/firm/predictor/firm_predictor.py}{53}{66}{Résidus par date sur les expositions.}\label{lst:rs:anatomy-of-a-predictor:residual}
\item \textbf{Résumer.} La moyenne de la série des IC, son ratio d'information et deux statistiques $t$, la seconde avec $h - 1$ retards de Newey--West
pour les cibles qui se chevauchent.
\omcode{code/firm/predictor/firm_predictor.py}{107}{118}{Le résumé des IC.}\label{lst:rs:anatomy-of-a-predictor:summary}
\item \textbf{Exécuter} \texttt{horizon\_ics()}, \texttt{study()} et \texttt{card()}~; puis \texttt{neutralisation} sur un marché simulé
avec \texttt{ind\_vol=0.30}, et \texttt{fig\_anatomy.py}.
\end{tutsteps}
\textbf{Pour aller plus loin.} Remplacer la \omterm{def:rs:anatomy-of-a-predictor:normalise}{transformation en rangs} par un z-score tronqué à 3, puis à 10, et observer l'écart type de l'IC~;
neutraliser aussi à l'exposition à la taille et vérifier que le retournement subsiste.
\end{tutorial}

\section{Construction~: la boîte à outils du prédicteur}\label{bld:rs:anatomy-of-a-predictor:predictor}

\begin{build}
\bfield{Rôle} Une seule implémentation des cibles, des normaliseurs, de la \omterm{def:rs:anatomy-of-a-predictor:neutralise}{neutralisation} et de l'IC pour chaque \omterm{def:rs:anatomy-of-a-predictor:predictor}{prédicteur} du livre, afin que
les chiffres issus de chapitres différents (et, plus tard, des livres~8 et~9) soient comparables.
\bfield{Interface} \texttt{forward\_return(ret, h)}, \texttt{excess(target, market)}, \texttt{residualise(y, X)} et
\texttt{neutralise}, \texttt{zscore(x, clip)}, \texttt{rank\_transform(x)}, \texttt{ic\_series(signal, target, method)},
\texttt{ic\_summary(ic, h, periods)}, \texttt{quantile\_spread(signal, target, q)}, \texttt{PredictorCard}.
\bfield{Règles} Les panels sont des (dates, noms) avec NaN pour les noms absents~; un \omterm{def:rs:anatomy-of-a-predictor:predictor}{prédicteur} à la ligne $t$ est connu à la clôture de $t$ et sa cible
commence en $t+1$~; les valeurs manquantes sont ignorées date par date, jamais imputées~; une date comptant moins de 20 noms n'a pas d'IC.
\bfield{Tests d'acceptation} \texttt{code/firm/predictor/tests/}~: capitalisation et trous dans les rendements à terme~; troncature et rangs~;
résidus orthogonaux aux expositions~; un IC de 0.05 retrouvé avec le bruit en $1/\sqrt N$~; un $t$ naïf supérieur de plus de 2.5 fois au $t$ HAC
pour un \omterm{def:rs:anatomy-of-a-predictor:predictor}{prédicteur} persistant et des cibles de 20 jours qui se chevauchent~; les champs d'une fiche.
\bfield{Pour aller plus loin} Les IC pondérés (par la liquidité)~; les IC par groupe (secteur, tranche de taille)~; un bootstrap de la moyenne de l'IC sur les dates.
\end{build}

\begin{omsources}
\item R.~C.~Grinold et R.~N.~Kahn, \emph{Active Portfolio Management}, 2e éd., McGraw-Hill, 2000.
\item E.~E.~Qian, R.~H.~Hua et E.~H.~Sorensen, \emph{Quantitative Equity Portfolio Management}, Chapman \& Hall/CRC, 2007.
\item W.~K.~Newey et K.~D.~West, ``A simple, positive semi-definite, heteroskedasticity and autocorrelation consistent covariance
matrix'', \emph{Econometrica} 55(3), 1987.
\item N.~Jegadeesh, ``Evidence of predictable behavior of security returns'', \emph{Journal of Finance} 45(3), 1990~; B.~N.~Lehmann,
``Fads, martingales, and market efficiency'', \emph{Quarterly Journal of Economics} 105(1), 1990.
\end{omsources}

\section{Exercices}

\begin{exercise}[$\star$]\label{exo:rs:anatomy-of-a-predictor:1}
Un univers compte 800 titres. Calculer l'écart type de l'IC d'une seule date pour un \omterm{def:rs:anatomy-of-a-predictor:predictor}{prédicteur} inutile, puis le nombre de dates indépendantes
nécessaires pour qu'un IC moyen de 0.01 atteigne $t = 2$.
\end{exercise}

\begin{exercise}[$\star$]\label{exo:rs:anatomy-of-a-predictor:2}
L'IC de rang quotidien d'un \omterm{def:rs:anatomy-of-a-predictor:predictor}{prédicteur} contre le rendement brut du lendemain vaut 0.02. Calculer son IC de rang contre le rendement
du lendemain en excès de l'indice équipondéré du marché, puis contre le rendement en excès du bêta propre de chaque action multiplié par le marché.
\end{exercise}

\begin{exercise}[$\star$]\label{exo:rs:anatomy-of-a-predictor:3}
Une série quotidienne d'IC a une moyenne de 0.012 et un écart type de 0.08. Calculer le \omterm{def:rs:anatomy-of-a-predictor:icir}{ratio d'information de l'IC} annualisé (ICIR).
\end{exercise}

\begin{exercise}[$\star\star$]\label{exo:rs:anatomy-of-a-predictor:4}
Les facteurs sectoriels expliquent 40~\% de la variance transversale des rendements mensuels, et un \omterm{def:rs:anatomy-of-a-predictor:predictor}{prédicteur} de sélection de titres sans
contenu sectoriel affiche un IC de 0.03 contre les rendements bruts. Calculer l'IC qu'il doit afficher contre les \omterm{def:rs:anatomy-of-a-predictor:residual}{rendements résiduels} sectoriels.
\end{exercise}

\begin{exercise}[$\star\star$]\label{exo:rs:anatomy-of-a-predictor:5}
Une cible de 20 jours est évaluée chaque jour pour un \omterm{def:rs:anatomy-of-a-predictor:predictor}{prédicteur} très persistant. D'après la proposition, indiquer de combien la $t$-statistique naïve
peut surestimer, puis calculer le $t$ auquel correspond une valeur naïve de 6.0.
\end{exercise}

\begin{exercise}[$\star\star$]\label{exo:rs:anatomy-of-a-predictor:6}
Donner une raison de préférer un IC de rang à un IC de Pearson et une raison de préférer l'IC de Pearson, en termes de ce que gagne un portefeuille construit à partir
du \omterm{def:rs:anatomy-of-a-predictor:predictor}{prédicteur}.
\end{exercise}

\begin{exercise}[$\star\star\star$]\label{exo:rs:anatomy-of-a-predictor:7}
\emph{Programmation.} Avec \texttt{ic\_series} et \texttt{ic\_summary}, mesurer le retournement à un jour contre le rendement du lendemain
séparément les jours où le rendement absolu du marché figure dans le premier décile et les autres jours. Indiquer quels jours portent le \omterm{def:rs:anatomy-of-a-predictor:predictor}{prédicteur}.
\end{exercise}

\begin{exercise}[$\star\star\star$]\label{exo:rs:anatomy-of-a-predictor:8}
\emph{Trouver la faille.} «~Notre \omterm{def:rs:anatomy-of-a-predictor:predictor}{prédicteur} à 60 jours affiche un IC moyen de 0.04 contre les rendements à terme de 60 jours, calculé chaque jour sur cinq ans~:
1\,260 dates, donc $t = 0.04/(0.10/\sqrt{1\,260}) = 14$.~»
\end{exercise}

\section{Problème~: deux IC pour un seul signal}

\begin{problem}[{Problème du week-end --- ce que les choix d'une mesure font aux chiffres d'un prédicteur}]\label{pb:rs:anatomy-of-a-predictor:1}
Le marché synthétique dans sa configuration par défaut (1\,000 titres, dix ans, graine 1), et une variante avec une volatilité sectorielle de
30~\% par an. Le \omterm{def:rs:anatomy-of-a-predictor:predictor}{prédicteur} est un retournement~: l'opposé d'un rendement passé.

\medskip\noindent\textbf{Partie I --- Horizon.}
\begin{enumerate}
\item Quel est l'IC de rang du retournement à un jour contre le rendement du lendemain, et sa statistique $t$~?
\item Et celui du retournement à cinq jours contre le rendement des cinq jours suivants~?
\item Quelles sont les deux paires croisées (un jour contre cinq, cinq jours contre un)~?
\item Laquelle des quatre est l'effet injecté, et pourquoi les autres sont-elles plus faibles~?
\item Quel est l'IC du retournement à un jour avec le rendement du deuxième jour suivant la décision~?
\end{enumerate}

\medskip\noindent\textbf{Partie II --- Cibles et \omterm{def:rs:anatomy-of-a-predictor:neutralise}{neutralisation}.}
\begin{enumerate}[resume]
\item Pourquoi la cible brute et la cible en excès du marché donnent-elles le même IC de rang~?
\item Quelle part de la variance des rendements à cinq jours expliquent le bêta et les secteurs sur le marché par défaut~?
\item Dans la variante à forte composante sectorielle, quelle est cette part, et quel rapport la proposition prédit-elle~?
\item Quels sont les trois IC de la \cref{fig:rs:anatomy-of-a-predictor:neutral} et leurs $t$-statistiques~?
\item Pourquoi neutraliser le \omterm{def:rs:anatomy-of-a-predictor:predictor}{prédicteur} s'ajoute-t-il à la \omterm{def:rs:anatomy-of-a-predictor:neutralise}{neutralisation} de la cible~?
\end{enumerate}

\medskip\noindent\textbf{Partie III --- Erreurs types.}
\begin{enumerate}[resume]
\item Quelles sont les $t$-statistiques naïve et corrigée du retournement à cinq jours contre des cibles à cinq jours~?
\item Pourquoi leur rapport est-il inférieur à $\sqrt5$~?
\item Que donne l'utilisation d'une date sur cinq, et pourquoi est-ce encore plus bas~?
\item Quel est l'écart type de l'IC d'une date avec 1\,000 titres et sans prévisibilité~?
\end{enumerate}

\medskip\noindent\textbf{Partie IV --- La fiche.}
\begin{enumerate}[resume]
\item Quels sont l'IC moyen, l'ICIR annualisé et la $t$-statistique de la fiche~?
\item Que dit la fiche de la demi-vie, et à partir de quels chiffres~?
\item Quel champ de la fiche changerait en premier si les fournisseurs de liquidité du marché devenaient plus nombreux~?
\item Énoncer le \emph{résultat nommé}~: les $t$-statistiques naïve et HAC du retournement à cinq jours pour des cibles qui se chevauchent, et le gain d'IC obtenu en neutralisant
la cible et le \omterm{def:rs:anatomy-of-a-predictor:predictor}{prédicteur} dans le marché à forte composante sectorielle.
\item Que doit dire un rapport de recherche lorsqu'il cite un IC~?
\item En une phrase~: qu'est-ce qu'un \omterm{def:rs:anatomy-of-a-predictor:predictor}{prédicteur}~?
\end{enumerate}
\end{problem}

\section{Questions d'entretien}

\begin{interviewq}[$\star$ \iqroles{chercheur}]\label{iq:rs:anatomy-of-a-predictor:1}
Qu'est-ce qu'un \omterm{def:rs:anatomy-of-a-predictor:ic}{coefficient d'information}, et quelles valeurs sont bonnes~?
\end{interviewq}

\begin{interviewq}[$\star\star$ \iqroles{chercheur, ingénieur ML}]\label{iq:rs:anatomy-of-a-predictor:2}
Vous calculez chaque jour l'IC d'un \omterm{def:rs:anatomy-of-a-predictor:predictor}{prédicteur} à 20 jours et obtenez $t = 8$. Y croyez-vous~?
\end{interviewq}

\begin{interviewq}[$\star\star$ \iqroles{chercheur}]\label{iq:rs:anatomy-of-a-predictor:3}
Pourquoi neutraliseriez-vous un \omterm{def:rs:anatomy-of-a-predictor:predictor}{prédicteur} par rapport aux secteurs~? Quand ne le feriez-vous pas~?
\end{interviewq}

\begin{interviewq}[$\star\star$ \iqroles{chercheur, trader}]\label{iq:rs:anatomy-of-a-predictor:4}
Un \omterm{def:rs:anatomy-of-a-predictor:predictor}{prédicteur} affiche un IC de 0.05 contre les rendements bruts et de 0.02 contre les \omterm{def:rs:anatomy-of-a-predictor:residual}{rendements résiduels}. Que se passe-t-il~?
\end{interviewq}

\begin{interviewq}[$\star\star$ \iqroles{ingénieur ML, chercheur}]\label{iq:rs:anatomy-of-a-predictor:5}
Rang IC ou IC de Pearson pour une prévision de rendement issue d'un apprentissage automatique~? Pourquoi~?
\end{interviewq}

\begin{interviewq}[$\star\star\star$ \iqroles{chercheur}]\label{iq:rs:anatomy-of-a-predictor:6}
Comment établiriez-vous l'écart type de l'IC d'une seule date sous l'absence de prévisibilité, et comment vous en serviriez-vous
pour dimensionner l'historique nécessaire à la détection d'un IC moyen de 0.02~?
\end{interviewq}
